% One-time imported draft. Editable native source; no Markdown regeneration.
\section{Direction and boundary cases to
predict}\label{007-bresenham-line-x-cases-direction-and-boundary-cases-to-predict}

{\def\LTcaptype{table} % do not increment counter
\begin{longtable}[]{@{}
  >{\RaggedRight\arraybackslash}p{(\linewidth - 4\tabcolsep) * \real{0.3333}}
  >{\RaggedRight\arraybackslash}p{(\linewidth - 4\tabcolsep) * \real{0.3333}}
  >{\RaggedRight\arraybackslash}p{(\linewidth - 4\tabcolsep) * \real{0.3333}}@{}}
\toprule\noalign{}
\begin{minipage}[b]{\linewidth}\RaggedRight
Endpoints before ordering
\end{minipage} & \begin{minipage}[b]{\linewidth}\RaggedRight
Case to inspect
\end{minipage} & \begin{minipage}[b]{\linewidth}\RaggedRight
Expected feature
\end{minipage} \\
\midrule\noalign{}
\endhead
\bottomrule\noalign{}
\endlastfoot
(0,2) to (7,2) & Horizontal & y stays two \\
(3,7) to (3,0) & Vertical and reversed & Canonical y runs zero through
seven \\
(1,1) to (3,7) & Steep increasing y & Some updates change only y \\
(1,6) to (6,2) & Negative slope & Successful y steps subtract one \\
(0,0) to (2,1) & Exact tie & Middle selected cell is (1,1) \\
(2,1) to (0,0) & Reversed exact tie & Same image as the previous row \\
(4,4) to (4,4) & Point & One selected cell \\
(0,7) to (7,0) & Diagonal at boundaries & Both axes change at each
update \\
\end{longtable}
}

These cases are not separate algorithms. Normalization, independent
decisions, and the post-store endpoint test handle them within one loop.
